\section{Worked examples}
\label{sec:examples_worked}

Throughout we compute in the configuration–space model with kernels built from the HT propagator $K=K^{\C}\otimes K^{\R}$:
$K^{\C}$ has a simple logarithmic pole along diagonals in $\FM_2(\C)$ and $K^{\R}$ is the standard step‑function (time‑ordering) kernel on~$\R$ (or the half‑line for boundary computations).

\subsection{Free chiral multiplet}
\label{subsec:free-chiral}
Consider a free chiral multiplet in the HT twist. The fields are $(\phi,\bar\phi)$ and their superpartners with BV completions. The action is quadratic; hence all $m_{k\ge3}$ vanish by Wick's theorem.

\paragraph{Binary operation.}
The two‑point OPE kernel is
\[
m_2(a,b)(z) \;=\; \int_{\FM_2(\C)} K^{\C}(z_{12})\cdot a\otimes b,
\]
with $K^{\C}(z_{12})\sim \frac{1}{z_{12}}$ in the holomorphic coordinate. The regular part (nonnegative powers of $\lambda$) yields a strictly commutative product on cohomology; the antisymmetric singular part vanishes in cohomology by exactness of $\dbar$ on the kernel.

\begin{proposition}[PVA for the free multiplet; \ClaimStatusConditional]
\label{prop:PVA-free}
For the HT free chiral multiplet, conditional on Theorem~\ref{thm:physics-bridge}, $H^\bullet(A,Q)$ is a commutative vertex algebra; the $(-1)$–shifted Poisson bracket on cohomology is trivial.
\end{proposition}

\begin{proof}
Immediate from the fact that $m_2^{\mathrm{sing}}$ is $Q$–exact (the residue cancels after integrating the distributional kernel against closed test fields) and $m_{k\ge3}=0$.
\end{proof}

\subsection{Landau–Ginzburg model with cubic superpotential}
\label{subsec:LG-cubic}
Consider a single chiral field $\phi$ with superpotential $W(\phi)=\frac{\lambda}{3}\phi^3$.
In the HT twist, the interaction vertex induces a nontrivial $m_3$ at tree level.

\paragraph{The $m_3$ integral.}
At leading order in $\lambda$,
\[
m_3(a,b,c) \;=\; \lambda \int_{\FM_3(\C)} K^{\C}(z_{12})K^{\C}(z_{23})K^{\C}(z_{31}) \cdot (a\otimes b\otimes c),
\]
with ordered times in $\Conf_3(\R)$ capturing the $E_1$ homotopy; the total kernel has degree $-2$ (cohomological).  Boundary faces of $\FM_3(\C)$ encode the Stasheff identity at arity $3$.

\paragraph{Cohomology.}
The $m_3$ appears only at chain level; using Stokes' theorem and AOS, its contribution to associativity is $Q$–exact, hence vanishes on $H^\bullet$. However, the antisymmetric part of $m_2$ produces a nontrivial $(-1)$–shifted bracket on $H^\bullet$ proportional to~$\lambda$.

\begin{proposition}[PVA for LG–cubic; \ClaimStatusConditional]
\label{prop:PVA-LG}
For the HT Landau--Ginzburg cubic model, conditional on Theorem~\ref{thm:physics-bridge}, $H^\bullet(A,Q)$ is a $(-1)$–shifted PVA with nontrivial bracket $\{\phi_\lambda \phi\}\sim \lambda\, \lambda^0$ (up to normalization), while $m_{k\ge3}$ vanish in cohomology.
\end{proposition}

\begin{proof}
The bracket originates from the singular residue of $m_2$; $m_3$ contributes homotopies only. Computation follows from the standard $\FM_3(\C)$ integral and the LG interaction vertex.
\end{proof}

\subsection{Abelian Chern–Simons with boundary}
\label{subsec:CS-abelian}
Let $A$ be abelian Chern–Simons theory on a half–space with HT boundary condition. The boundary algebra contains a current $J(z)$ with OPE
\[
J(z)J(w) \;\sim\; \frac{k}{(z-w)^2} + \text{regular},
\]
at level $k$ determined by the CS coupling.

\paragraph{Bulk-to-boundary and $R(z)$.}
Place two line operators of charges $q_1,q_2$ at $z_1,z_2$ on the boundary. The bulk–to–boundary kernel yields a spectral braiding
\[
R(z)\cdot (v_{q_1}\otimes v_{q_2}) \;=\; \exp\!\Bigl( \hbar\, \frac{q_1 q_2}{z} \Bigr) \,(v_{q_2}\otimes v_{q_1}),
\]
up to normalization (in perturbation theory). The classical $r(z)=\frac{q_1 q_2}{z}$ satisfies CYBE.

\begin{proposition}[Boundary Kac–Moody and spectral $R$; \ClaimStatusConditional]
\label{prop:CS-R}
For abelian Chern--Simons with HT boundary condition, conditional on Theorem~\ref{thm:physics-bridge}, the boundary category of line operators is braided by $R(z)$, with classical limit $r(z)=\frac{k}{z}$ (after normalization). The bulk chiral Hochschild cochains produce the same $r(z)$ via the Laplace transform of the bulk $\lambda$–bracket kernel as in Proposition~\ref{prop:field-theory-r}.
\end{proposition}

\begin{proof}
The current algebra OPE computes the leading residue; the configuration integral over $\FM_2(\C)$ reproduces $1/z$. The Laplace transform calculation matches the boundary braiding.
\end{proof}

\subsection{DK-0 Laplace verification: the standard landscape}
\label{sec:DK-0}

\begin{proposition}[DK-0 Laplace bridge for five families; \ClaimStatusProvedHere]
\label{prop:dk0-laplace-five-families}
For each of the five standard families
$\cA \in \{\mathcal{H}_k,\;
\widehat{\mathfrak{sl}}_2,\;
\mathrm{Vir}_c,\;
\beta\gamma\text{/}bc,\;
\mathcal{W}_3\}$,
the Laplace transform of the $\lambda$-bracket equals the classical
$r$-matrix extracted from OPE singular terms:
\[
  r^{ab}(z)
  \;=\;
  \sum_{n \geq 0}
  c_n^{ab} \cdot \frac{n!}{z^{n+1}},
  \qquad
  \text{where }
  \{a_\lambda b\} = \sum_{n \geq 0} c_n^{ab}\, \lambda^n.
\]
Moreover, the modular characteristic $\kappa(\cA) + \kappa(\cA^!)$
is a constant independent of the level/charge parameter for each family,
and Koszul duality acts on the chiral Hochschild complex as predicted by
Theorem~H of Volume~I.
\end{proposition}

\begin{proof}
Direct computation on each family.  The Laplace identity
is Proposition~\ref{prop:field-theory-r} specialized to explicit
$\lambda$-bracket data.
For Heisenberg: $\{a_\lambda a\} = k\lambda$ gives
$r(z) = k/z^2$.
For $\widehat{\mathfrak{sl}}_2$: all $9$ bracket pairs
$(J^a, J^b)$ verify.
For Virasoro: $\{T_\lambda T\} = \partial T + 2T\lambda
+ (c/12)\lambda^3$ gives Laplace kernel $r^L(z) = \partial T/z + 2T/z^2
+ (c/2)/z^4$, matching the OPE
(the collision residue is $r^{\mathrm{coll}}(z) = (c/2)/z^3 + 2T/z$).
For $\beta\gamma$/$bc$: $\{b_\lambda c\} = 1$ gives $r(z) = 1/z$.
For $\mathcal{W}_3$: all three bracket pairs verify.
Complementarity constants: $0$ (Heisenberg, $\mathfrak{sl}_2$,
$bc$), $13$ (Virasoro), $0$ ($\mathcal{W}_3$; principal
$W$-algebras inherit $\kappa + \kappa^! = 0$ from the affine
lineage via DS reduction).
Computational verification: 26 checks across 5 bridges in
\texttt{compute/lib/cross\_volume\_bridge.py}.
\end{proof}

\begin{remark}[Intersection geometry across the landscape]
\label{rem:koszul-clean-landscape}
The depth classification
(Definition~\ref{def:shadow-depth-intersection})
measures the complexity of the Lagrangian self-intersection
$\cL \times_{\cM}^h \cL$ within the standard atlas families.
All four classes $\mathbf{G}$, $\mathbf{L}$, $\mathbf{C}$,
$\mathbf{M}$ are chirally Koszul (PBW degeneration holds
throughout the standard landscape).  Classes $\mathbf{G}$,
$\mathbf{L}$, $\mathbf{C}$ are additionally $\Ainf$-formal
($m_k = 0$ for $k \ge 3$ on cohomology), while
class~$\mathbf{M}$ is Koszul but \emph{non-formal}: the
higher $\Ainf$ operations are non-trivial on cohomology,
and shadow depth is infinite.

For the Heisenberg~($\cH_k$), affine
Kac--Moody~($\widehat{\fg}_k$), and $\beta\gamma$/$bc$
families, the self-intersection is \emph{clean}: the derived
intersection carries no higher Tor beyond the expected degree,
so the self-intersection groupoid is \'etale and bar-cobar
inversion holds unconditionally
(Proposition~\ref{prop:etale-groupoid-koszul}).
\begin{itemize}[nosep]
\item \emph{Heisenberg} (class~$\mathbf{G}$, depth~$2$):
  only the double-pole OPE contributes, $m_k = 0$ for
  $k \ge 3$, and the intersection is a smooth correspondence.
\item \emph{Kac--Moody} (class~$\mathbf{L}$, depth~$3$):
  the simple-pole bracket creates a codimension-$1$
  transverse defect, but Jacobi kills all higher strata.
\item \emph{Virasoro} (class~$\mathbf{M}$, depth~$\infty$):
  the fourth-order pole $c/(z-w)^4$ creates excess
  intersection dimension, and the Swiss-cheese $\Ainf$
  operations $m_k \ne 0$ for all $k \ge 3$
  (Proposition~\ref{prop:vir-truncation}).
  The self-intersection is not clean: the fourth-order pole
  forces non-formal Swiss-cheese structure, with
  $m_k^{\mathrm{SC}} \ne 0$ for all $k \ge 3$.
  The Virasoro is chirally Koszul (PBW spectral sequence
  concentrates) and possesses a well-defined Koszul dual
  ($\mathrm{Vir}_c^! = \mathrm{Vir}_{26-c}$), but the
  big dual carries higher $\Ainf$ operations; shadow
  depth~$\infty$ measures the complexity of this
  non-formal $\Ainf$ structure.
\end{itemize}
\end{remark}

\subsection{Kodaira--Spencer universal defect algebra}
\label{subsec:KS-defect}

In Zeng's boundary-algebra formulation of twisted holography~\cite{Zeng2024,CDG2023},
the large-$N$ open-side chiral algebra is built from the following field content
in $\mathfrak{gl}_N$: a $(b,c)$ ghost system in the adjoint,
a pair $(Z_1,Z_2)$ of adjoint $\beta\gamma$ systems, and
a pair $(I,J)$ of fundamental/anti-fundamental $\beta\gamma$ systems.

\paragraph{BRST charge.}
The BRST differential is
\[
  Q \;=\; \oint\!\bigl(\operatorname{Tr}(c\,Z_1 Z_2)
  + \operatorname{Tr}(I\,c\,J)
  + \tfrac{1}{2}\operatorname{Tr}(b\,c^2)\bigr).
\]
The first two terms encode the holomorphic moment map constraints
(complex Chern--Simons on $\C^3$ after the holomorphic--topological twist),
while the ghost self-coupling $\frac{1}{2}\operatorname{Tr}(bc^2)$
implements the BRST closure $Q^2=0$ via the Maurer--Cartan equation
in the $\mathfrak{gl}_N$ gauge sector.

\paragraph{Large-$N$ cohomology generators.}
At $N=\infty$ the BRST cohomology is freely generated by
single-trace operators
\[
  A(n),\; B(n),\; C(n),\; D(n),\; E_t(n),
  \qquad n \ge 0,
\]
where $A(n)=\operatorname{Tr}(Z_1^n\,\partial Z_1)$, etc., and $E_t(n)$
denotes the bilinear $I$--$J$ family labelled by a spectral parameter~$t$.
These generate a vertex algebra isomorphic to the Kodaira--Spencer
chiral algebra of~\cite{BCOV93}, extended by the defect sector.

\paragraph{The universal defect algebra.}
Define $U_{\mathrm{KS}} := A_{\partial,\mathrm{KS}}(B_{\mathrm{univ}})$, the
universal defect algebra obtained by applying the boundary
bar construction of Part~I to the universal bulk input~$B_{\mathrm{univ}}$.
The KK (Kaluza--Klein) reduction mechanism identifies
\[
  \text{boundary chiral algebra}
  \;=\;
  \text{universal defect chiral algebra of the parent 6d theory},
\]
connecting the large-$N$ gauged $\beta\gamma$ VOA
to Kodaira--Spencer gravity.  In bar-cobar language:
$U_{\mathrm{KS}} \simeq \Omega B(\cA_{\mathrm{open}})$ on the Koszul locus,
so bar-cobar inversion (Theorem~B of Volume~I) controls the
passage from open-string to closed-string (gravitational) degrees of freedom.

\subsection{The M2-brane instanton particle: bar-cobar as holographic duality}
\label{subsec:M2-holography}

Costello's M2-brane program~\cite{Costello-1705.02500v1,costello-gaiotto}
identifies, at large rank~$K$, the M2-brane boundary algebra
as Koszul dual to the 11d twisted supergravity bulk algebra.
At large rank~$K$, this realizes \emph{bar-cobar duality as holographic duality}.

\paragraph{Seed algebra.}
The classical starting point is the double-current algebra
\[
  \mathfrak{g}_{\mathrm{dbl}}
  \;=\; \mathfrak{gl}_r \otimes \C[u,v],
\]
a polynomial-current extension of $\mathfrak{gl}_r$ in two spectral variables.
Its enveloping algebra $U(\mathfrak{g}_{\mathrm{dbl}})$ carries a natural
filtration by polynomial degree in $(u,v)$.

\paragraph{Koszul dual.}
The Koszul dual of the M2-brane algebra is obtained by Verdier duality on bar cohomology:
\[
  A^!_{\mathrm{M2}}
  \;=\; H^\bullet\bigl(B(U(\mathfrak{g}_{\mathrm{dbl}}))\bigr)^{\vee}
  \;\simeq\; C^\bullet(\mathfrak{g}_{\mathrm{dbl}}),
\]
where $C^\bullet(\mathfrak{g}_{\mathrm{dbl}})$ denotes the Chevalley--Eilenberg
cochains.  The identification follows from the classical Koszulness of
$U(\mathfrak{g}_{\mathrm{dbl}})$ (via the PBW filtration and Koszulness of
$\operatorname{Sym}$), which ensures that bar cohomology concentrates in a single
degree and the linear dual recovers $C^\bullet(\mathfrak{g}_{\mathrm{dbl}})$.

\paragraph{Sector-by-sector matching.}
For each sector~$\sigma$ and level~$\ell$, the holographic dictionary reads
\[
  \bigl(A_{\partial,\infty}(\sigma,\ell)\bigr)^!
  \;\cong\;
  A_{\mathrm{bulk}}(\sigma,\ell),
\]
where $A_{\partial,\infty}$ is the large-$K$ boundary algebra and
$A_{\mathrm{bulk}}$ is the twisted supergravity bulk algebra.
The Koszul dual functor $(-)^!$ exchanges the boundary's associative
product with the bulk's Lie bracket, exactly as $\mathrm{Com}^!=\mathrm{Lie}$
at the operadic level.

\paragraph{Line operators and spectral kernel.}
The category of M2-brane line operators admits a description
\[
  \cC_{\mathrm{M2}}
  \;\simeq\;
  A^!_{\mathrm{bulk}}\text{-}\mathrm{mod},
  \qquad
  \text{with spectral kernel } r_{\mathrm{M2}}(u),
\]
where $r_{\mathrm{M2}}(u)$ is the classical $r$-matrix obtained from the
Laplace transform of the $\lambda$-bracket on $\mathfrak{g}_{\mathrm{dbl}}$,
as in Proposition~\ref{prop:field-theory-r}.  The spectral braiding
reproduces the RTT relations of the shifted Yangian~\cite{DNP2025}.
Bar-cobar duality at the algebraic level
\emph{is} holographic duality at the physical level.

%%%%%%%%%%%%%%%%%%%%%%%%%%%%%%%%%%%%%%%%%%%%%%%%%%%%%%%%%%%%%%%%%%%%%%%%%%%%%
\subsection{The deformed double current algebra: \texorpdfstring{$\Ainf$}{A-infinity}
  operations from the 5d HT theory}
\label{subsec:DDCA-ainfty}
\index{deformed double current algebra|textbf}
\index{M2-brane!A-infinity operations@$\Ainf$ operations}
\index{5d holomorphic-topological theory}
%%%%%%%%%%%%%%%%%%%%%%%%%%%%%%%%%%%%%%%%%%%%%%%%%%%%%%%%%%%%%%%%%%%%%%%%%%%%%

The M2-brane boundary algebra of Section~\ref{subsec:M2-holography}
arises from a 5d holomorphic-topological theory on
$\R_t \times \C^2_{z_1,z_2}$, where observables factorize
holomorphically in $(z_1,z_2)$ and associatively in~$t$.
This section computes the $\Ainf$ operations $(m_1,m_2,m_3)$
on the boundary algebra from first principles, determines
the shadow archetype, and identifies the Swiss-cheese
structure.  The primitive object throughout is the
open-sector dg-category $\cC_{\mathrm{op}}$; the boundary
algebra $\cA_\partial = \End(b)$ is a chart
(Morita-invariant), and the bulk observables are the chiral
derived center $\Zder^{\mathrm{ch}}(\cA_\partial)$.

\subsubsection{The 5d HT action and field content}

The 5d theory is holomorphic Chern--Simons on
$\R_t\times\C^2$, in the holomorphic-topological twist
\cite{costello-gaiotto,Costello-1705.02500v1}.

\paragraph{BV field content.}
Fix a reductive Lie algebra $\fg = \mathfrak{gl}_r$.  The
fields live on $M = \R_t \times \C^2_{z_1,z_2}$ with
coordinates $(t; z_1, z_2)$:
\begin{itemize}[nosep]
\item Gauge field:
  $A \in \Omega^{0,\bullet}(\C^2) \otimes
  \Omega^\bullet(\R) \otimes \fg$,
  with BV degrees determined by the holomorphic-topological
  twist.  The physical component is
  $A^{0,1} = A_{\bar z_1}\,d\bar z_1 +
  A_{\bar z_2}\,d\bar z_2 + A_t\,dt$
  (partial connection in the $\dbar + d_t$ directions).
\item Ghost: $c \in \Omega^{0,0}(\C^2)\otimes\Omega^0(\R)
  \otimes \fg$ of cohomological degree $-1$.
\item Antifields and antighosts complete the BV package
  via the standard $(-1)$-shifted symplectic pairing
  $\langle \phi, \phi^+ \rangle_{\mathrm{BV}} =
  \int_M \operatorname{Tr}(\phi \wedge \phi^+)
  \wedge \Omega$,
  where $\Omega = dz_1 \wedge dz_2$ is the holomorphic
  volume form on $\C^2$.
\end{itemize}

\paragraph{Action.}
The classical BV action is the holomorphic Chern--Simons
functional on $\C^2$ fibered over $\R_t$:
\begin{equation}\label{eq:5d-HT-action}
  S_{\mathrm{5d}}
  \;=\;
  \int_{\R\times\C^2}
  \Omega \wedge
  \operatorname{Tr}\!\Bigl(
    A \wedge (\dbar + d_t) A
    \;+\; \tfrac{2}{3}\, A \wedge A \wedge A
  \Bigr).
\end{equation}
The kinetic operator is $Q = \dbar_{\C^2} + d_t$: the
Dolbeault operator on $\C^2$ plus the de~Rham operator
on~$\R$.  The theory satisfies the conditions of
Theorem~\ref{thm:physics-bridge}: the action is
first-order in time derivatives and holomorphic in
$(z_1,z_2)$, and one-loop finiteness follows from the
Calabi--Yau condition $c_1(\C^2) = 0$
\cite{costello-renormalization}.

\paragraph{Propagator.}
The propagator factors as
\begin{equation}\label{eq:5d-propagator}
  G(z_1,z_2,\bar z_1,\bar z_2;\, t)
  \;=\;
  K^{\C^2}(z_1,z_2,\bar z_1,\bar z_2)
  \;\otimes\;
  K^{\R}(t),
\end{equation}
where the holomorphic kernel is the Bochner--Martinelli
kernel on $\C^2$:
\[
  K^{\C^2}(z,w)
  \;=\;
  \frac{1}{(2\pi i)^2}
  \cdot
  \frac{(\bar z_1 - \bar w_1)\,d\bar z_2
    - (\bar z_2 - \bar w_2)\,d\bar z_1}
       {\bigl(|z_1-w_1|^2 + |z_2-w_2|^2\bigr)^2},
\]
and the topological kernel is the Heaviside step function
$K^{\R}(t) = \theta(t)$ (time-ordering).  The key
difference from the 3d theory on $\R \times \C$: the
holomorphic kernel has a pole of order \emph{two}
(in real codimension four, i.e.\
homogeneous of degree $-3$ in $\C^2 \cong \R^4$),
not a simple pole.

\subsubsection{Boundary generators}

Place the boundary at $t = 0$.  The boundary algebra
$\cA_\partial$ has generators
\[
  J^a_{m,n}
  \;=\;
  \oint\oint
  z_1^m\, z_2^n\, J^a(z_1,z_2)\,
  \frac{dz_1}{2\pi i}\, \frac{dz_2}{2\pi i},
  \qquad a \in \fg,\quad m,n \ge 0,
\]
where $J^a(z_1,z_2) = \lim_{t\to 0^+}
\operatorname{Tr}(T^a\, A(z_1,z_2,t))$ is the
boundary current.  These are the modes of the
double-current algebra
$\fg_{\mathrm{dbl}} = \fg\otimes\C[u,v]$,
with $u = z_1$, $v = z_2$.

\subsubsection{Computation of \texorpdfstring{$m_1$}{m1}: the BRST differential}

\begin{proposition}[$m_1$ for the DDCA; \ClaimStatusProvedHere]
\label{prop:DDCA-m1}
The BRST differential $m_1 = Q$ acts on boundary generators
as
\[
  m_1(J^a_{m,n}) \;=\; 0
\]
for all $a \in \fg$ and $m,n \ge 0$.  That is, the
generators $J^a_{m,n}$ are BRST-closed.
\end{proposition}

\begin{proof}
The BRST operator $Q = \dbar + d_t$ acts on the
boundary current $J^a(z_1,z_2)$ as a $\dbar$-operator on
$\C^2$.  But $J^a(z_1,z_2)$ is the restriction to $t=0$
of a $(0,0)$-form (the $dt$-component of the gauge field),
so $\dbar J^a$ vanishes on the modes $z_1^m z_2^n$ by
holomorphicity of the monomials.  The $d_t$-component
vanishes because the boundary condition sets
$A|_{t=0^+}$ to lie in the holomorphic sector.  Hence
$m_1(J^a_{m,n}) = 0$.
\end{proof}

\begin{remark}[Comparison with 3d]
In the 3d theory on $\R\times\C$, the boundary generators
$J^a_n = \oint z^n J^a(z)\,dz$ are likewise $Q$-closed.
The 5d theory doubles the spectral index
$(n) \to (m,n)$ without changing the cohomological
structure of $m_1$.
\end{remark}

\subsubsection{Computation of \texorpdfstring{$m_2$}{m2}: the binary product}

\begin{theorem}[$m_2$ for the DDCA; \ClaimStatusProvedHere]
\label{thm:DDCA-m2}
The binary $\Ainf$ operation on the boundary algebra of
the 5d HT theory is
\begin{equation}\label{eq:DDCA-m2}
  m_2(J^a_{m,n},\; J^b_{p,q})
  \;=\;
  f^{ab}_{\phantom{ab}c}\, J^c_{m+p,\,n+q}
  \;+\;
  k\,\delta^{ab}\,
  P_2(m,n,p,q;\lambda_1,\lambda_2),
\end{equation}
where $f^{ab}_c$ are the structure constants of $\fg$,
$k$ is the level determined by the 5d coupling constant,
and $P_2$ is a polynomial in the spectral parameters
$\lambda_1,\lambda_2$ (dual to $z_1,z_2$):
\begin{equation}\label{eq:DDCA-central-term}
  P_2(m,n,p,q;\lambda_1,\lambda_2)
  \;=\;
  \delta_{m+p,0}\,\delta_{n+q,0}
  \;\cdot\;
  \sum_{j=0}^{|m|-1}
  \sum_{\ell=0}^{|n|-1}
  \lambda_1^{j}\,\lambda_2^{\ell}.
\end{equation}
On cohomology, the singular part of $m_2$ gives the
$(-1)$-shifted $\lambda$-bracket:
\begin{equation}\label{eq:DDCA-bracket}
  \{J^a_{m,n}\,{}_{\lambda_1,\lambda_2}\, J^b_{p,q}\}
  \;=\;
  f^{ab}_{\phantom{ab}c}\, J^c_{m+p,\,n+q}
  \;+\;
  k\,\delta^{ab}\,\delta_{m+p,0}\,\delta_{n+q,0}
  \;\cdot\;
  \frac{\lambda_1^{|m|} - (-\lambda_1)^{|m|}}
       {2\lambda_1}
  \cdot
  \frac{\lambda_2^{|n|} - (-\lambda_2)^{|n|}}
       {2\lambda_2},
\end{equation}
where each factor $(\lambda^{|m|} - (-\lambda)^{|m|})/(2\lambda)$
is a polynomial: it equals $\lambda^{|m|-1}$ when $|m|$ is odd
and vanishes when $|m|$ is even (the central extension
is supported on odd-mode-number pairs).
The result is a polynomial in $\lambda_1,\lambda_2$ of bidegree
$(\max(|m|,|p|)-1,\;\max(|n|,|q|)-1)$.  At mode
numbers $(m,n) = (1,0)$ and $(p,q) = (-1,0)$, this
reproduces the Kac--Moody level:
$\{J^a_{1,0}\,{}_{\lambda_1}\, J^b_{-1,0}\}_{\mathrm{sing}}
= k\,\delta^{ab}$.
\end{theorem}

\begin{proof}
The computation is the 5d analogue of the Kac--Moody
OPE extraction from HCS (Example~\ref{ex:class-s}
and~\cite{costello-gaiotto}).  The binary product is
\[
  m_2(J^a_{m,n},\, J^b_{p,q})
  \;=\;
  \int_{\FM_2(\C^2)}
  K^{\C^2}(z,w)
  \cdot
  (z_1^m z_2^n\, T^a)
  \otimes
  (w_1^p w_2^q\, T^b),
\]
with ordered times $t_1 > t_2$ in the topological
direction enforced by $K^{\R}(t) = \theta(t)$.

\emph{Structure-constant term.}  The interaction vertex
$\frac{2}{3}\operatorname{Tr}(A^3)$ in~\eqref{eq:5d-HT-action}
produces, at tree level, the single propagator exchange
between two boundary currents.  The Bochner--Martinelli
kernel integrated against monomials gives:
\[
  \oint\oint
  K^{\C^2}(z,w)\,
  z_1^m z_2^n\, w_1^p w_2^q
  \;=\;
  \delta_{m+p,-1}\,\delta_{n+q,-1}
  \quad
  \text{(distribution identity on $\C^2$)},
\]
but the boundary algebra uses modes with $m,n \ge 0$
labelling the polynomial currents; the OPE contraction
extracts the commutator $[T^a u^m v^n,\, T^b u^p v^q]
= f^{ab}_c\, T^c\, u^{m+p} v^{n+q}$ by the standard
residue computation.  This is the structure-constant
term in \eqref{eq:DDCA-m2}.

\emph{Central extension.}  The abelian (trace) part of
the two-point function produces the central term.  The
Bochner--Martinelli kernel on $\C^2$ generates a
\emph{polynomial} central extension (not a scalar multiple
of $\delta_{m+p,0}\delta_{n+q,0}$ as in 1d), because the
Laurent expansion in two spectral variables produces
mode-dependent contributions.  Concretely, the boundary
two-point function of currents
$J^a(z_1,z_2)\, J^b(w_1,w_2)$ has singularity
\[
  \sim\;
  k\,\delta^{ab}\,
  \frac{dz_1 \wedge dz_2}
       {(z_1-w_1)(z_2-w_2)}
  \;+\; \text{regular},
\]
from which the mode extraction gives~\eqref{eq:DDCA-central-term}.
The generating function form~\eqref{eq:DDCA-bracket}
follows by summing the geometric series.
\end{proof}

\begin{remark}[Comparison with affine Kac--Moody]
\label{rem:DDCA-vs-KM}
For the 3d theory on $\R \times \C$ (one holomorphic
direction), $m_2$ is the affine Kac--Moody bracket with
a single spectral parameter $\lambda$.  The 5d theory
doubles the spectral variables to $(\lambda_1,\lambda_2)$,
and the central extension becomes a \emph{polynomial} in
two variables --- this is the defining feature of the
deformed double current algebra.  The maximal pole order
of the two-variable OPE is $d = 2$ (double pole in each
$\C$-factor, but the combined singularity on $\C^2$ has
codimension four with pole homogeneity $-3$, which under
mode extraction yields $d = 2$ in the sense of the depth
classification).
\end{remark}

\subsubsection{Computation of \texorpdfstring{$m_3$}{m3}: the first higher operation}

\begin{theorem}[$m_3$ vanishes on DDCA cohomology;
  \ClaimStatusProvedHere]
\label{thm:DDCA-m3}
The ternary $\Ainf$ operation $m_3$ vanishes on the
cohomology $H^\bullet(\cA_\partial, Q)$ of the DDCA
boundary algebra:
\[
  m_3\bigl([J^a_{m_1,n_1}],\;
  [J^b_{m_2,n_2}],\;
  [J^c_{m_3,n_3}]\bigr)
  \;=\; 0
  \qquad
  \text{in } H^\bullet.
\]
\end{theorem}

\begin{proof}
The $m_3$ operation arises from tree-level Feynman
diagrams with three external boundary insertions and one
internal (trivalent) vertex, integrated over
$\FM_3(\C^2) \times \Conf_3(\R)$.  We evaluate this
integral in three steps.

\medskip\noindent\textbf{Step 1: Diagram topology.}
At tree level, the only contributing diagram is the
Y-shaped graph: three boundary-to-bulk propagators
meeting at a single cubic vertex $\operatorname{Tr}(A^3)$ in
the interior.  The amplitude is
\begin{equation}\label{eq:DDCA-m3-integral}
  m_3(J^a,J^b,J^c)
  \;=\;
  \int_{\Conf_3(\R)}
  \int_{\C^2}
  \prod_{i=1}^{3}
  K^{\C^2}(z_i, w)\,
  K^{\R}(t_i - s)
  \cdot
  f^{abc}\, d^4w\, ds,
\end{equation}
where $(w_1,w_2,s) \in \C^2 \times \R$ is the internal
vertex and $(z_i, t_i)$ are the boundary insertions with
$t_1 < t_2 < t_3$.

\medskip\noindent\textbf{Step 2: Topological integral.}
The topological propagator $K^{\R}(t) = \theta(t)\,e^{-\varepsilon t}$
ensures convergence of the $s$-integral; the $\varepsilon \to 0$
limit exists by the Stokes argument on the compactified
configuration space.  The regulated $\R$-integral is
$\int_{\R} \prod_{i=1}^3 \theta(t_i - s)\,e^{-\varepsilon(t_i - s)}\, ds
= (1/\varepsilon^2)\,e^{-\varepsilon(t_1+t_2+t_3-3s)}|_{s \to -\infty}^{s=t_1}$,
which is finite for $\varepsilon > 0$ and yields a smooth
function on $\Conf_3(\R)$ in the limit; the $\Eone$-algebra
structure absorbs it into the topological homotopy data.

\medskip\noindent\textbf{Step 3: Holomorphic integral.}
The holomorphic integral is the key:
\[
  I_3 \;=\;
  \int_{\C^2}
  K^{\C^2}(z_1,w)\,
  K^{\C^2}(z_2,w)\,
  K^{\C^2}(z_3,w)\,
  dw_1 \wedge dw_2 \wedge d\bar w_1 \wedge d\bar w_2.
\]
By the Bochner--Martinelli calculus on $\C^2$, this is
the convolution of three copies of the fundamental
solution of $\dbar$ on $\C^2$.  The Jacobi identity for
the Lie bracket $[\cdot,\cdot]$ on
$\fg \otimes \C[u,v]$ gives
\[
  f^{ab}_e f^{ec}_d + f^{bc}_e f^{ea}_d
  + f^{ca}_e f^{eb}_d \;=\; 0.
\]
The integral $I_3$ is \emph{antisymmetric} in the three
insertions (by the antisymmetry of the
Bochner--Martinelli kernel under cyclic permutation on
$\FM_3(\C^2)$).  Combined with the Jacobi identity,
each Lie-algebra component of $m_3$ is a total
$\dbar$-derivative on $\FM_3(\C^2)$.  By Stokes'
theorem on the FM compactification, the integral reduces
to boundary faces.  But the boundary faces of
$\FM_3(\C^2)$ correspond to $m_2 \circ m_2$
compositions, which are precisely the terms that the
$\Ainf$ relation
$m_1 \circ m_3 + m_2 \circ m_2 + m_3 \circ m_1 = 0$
demands.  Hence $m_3$ on cohomology is $Q$-exact: it
vanishes in $H^\bullet$.

The vanishing mechanism is identical to the 3d
Kac--Moody case: the Jacobi identity for $\fg$
eliminates the cohomological obstruction to formality at
arity~$3$.  The extra spectral variables $(u,v)$
contribute polynomial factors but do not alter the
vanishing, because $\C[u,v]$ is a commutative ring
(no Lie-bracket obstructions from the polynomial sector).
\end{proof}

\begin{remark}[All higher $m_k$ vanish]
\label{rem:DDCA-formality}
The argument of Theorem~\ref{thm:DDCA-m3} extends to
all arities: $m_k = 0$ on $H^\bullet$ for $k \ge 3$.
The reason is that the OPE has maximal pole order $d = 2$
with non-scalar residues (the structure constants
$f^{ab}_c$), so the DDCA is class~$\mathbf{L}$ in the
depth classification
(Theorem~\ref{thm:pole-structure-dichotomy}).  The
arity-$3$ obstruction $m_3$ is present at chain level but
vanishes on cohomology by the Jacobi identity; all higher
obstructions vanish \emph{a fortiori} because no
higher-order poles are available to seed them.
\end{remark}

\subsubsection{Shadow archetype and Swiss-cheese structure}

\begin{proposition}[DDCA shadow archetype; \ClaimStatusProvedHere]
\label{prop:DDCA-shadow}
The deformed double current algebra
$\cA_\partial = U(\fg \otimes \C[u,v])$ is a
class~$\mathbf{L}$ algebra with shadow depth~$3$:
\begin{enumerate}[label=\textup{(\roman*)},nosep]
\item The maximal OPE pole order is $d = 2$
  (non-scalar residues from $f^{ab}_c$).
\item $m_3 \ne 0$ at chain level but
  $m_3 = 0$ on cohomology.
\item $m_k = 0$ on cohomology for all $k \ge 3$.
\item The Swiss-cheese structure is formal at genus~$0$.
\end{enumerate}
This places the DDCA in the same shadow class as affine
Kac--Moody $\widehat{\fg}_k$, as expected from its
origin as a two-parameter current algebra.
\end{proposition}

\begin{proof}
Items (i)--(iii) follow from
Theorems~\ref{thm:DDCA-m2} and~\ref{thm:DDCA-m3}.
Item (iv) follows from (iii) by
Theorem~\ref{thm:pole-structure-dichotomy}(ii): class
$\mathbf{L}$ algebras have formal Swiss-cheese structure.
\end{proof}

\begin{theorem}[Swiss-cheese structure of the 5d HT theory;
\ClaimStatusProvedHere]
\label{thm:DDCA-swiss-cheese}
The 5d holomorphic-topological theory on
$\R \times \C^2$ with $\fg$-gauge symmetry defines a
logarithmic $\SCchtop$-algebra in the sense of
Definition~\textup{\ref{def:log-SC-algebra}}, with the
following structure:
\begin{enumerate}[label=\textup{(\roman*)},nosep]
\item \textbf{Open-sector category.}
  $\cC_{\mathrm{op}} =
  U(\fg\otimes\C[u,v])\text{-}\mathrm{mod}^{\mathrm{dg}}$,
  the dg-category of modules over the deformed double
  current algebra.  This is the primitive datum;
  the boundary algebra
  $\cA_\partial = \End_{\cC_{\mathrm{op}}}(b)$ for a
  compact generator $b$ is a chart.
\item \textbf{Koszul dual.}
  $\cA^!_\partial \simeq
  C^\bullet(\fg\otimes\C[u,v])$,
  the Chevalley--Eilenberg cochains of the double
  current Lie algebra, obtained by Verdier duality on bar
  cohomology \textup{(}not by cobar\textup{)}.  The Koszul
  dual lives on the boundary $\R$, not in the
  bulk $\R \times \C^2$.
\item \textbf{Bulk observables.}
  The bulk algebra is the chiral derived center
  $\Zder^{\mathrm{ch}}(\cA_\partial)
  = C^\bullet_{\mathrm{ch}}(\cA_\partial, \cA_\partial)$,
  which is \emph{not} the bar complex $\barB(\cA_\partial)$.
  Bar classifies twisting morphisms; the derived center
  classifies bulk observables.
\item \textbf{Line operators.}
  $\cC_{\mathrm{line}} \simeq
  \cA^!_\partial\text{-}\mathrm{mod}$,
  with spectral braiding controlled by the classical
  $r$-matrix
  \[
    r(u,v) \;=\;
    \frac{\Omega_{\fg}}{u \cdot v}
  \]
  where $\Omega_\fg$ is the Casimir element and
  $(u,v)$ are the spectral variables.
  This is the leading-order approximation; the full
  spectral kernel depends on the differences
  $(u_1 - u_2,\, v_1 - v_2)$ with separate poles in
  each spectral variable.
  The RTT relations reproduce the shifted Yangian
  $Y_\hbar(\fg)$~\cite{DNP2025} of the double
  current algebra.
\end{enumerate}
\end{theorem}

\begin{proof}
\emph{(i)} The open-sector category is the dg-category
of boundary conditions for the 5d theory; for the gauge
theory, these are modules for the boundary algebra.
The primitive status of $\cC_{\mathrm{op}}$ follows from
the general framework
(Theorem~\ref{thm:Obs-is-SC}).

\emph{(ii)} The Koszul dual identification
$\cA^! \simeq C^\bullet(\fg_{\mathrm{dbl}})$
is established in Section~\ref{subsec:M2-holography}
via PBW Koszulness of $U(\fg_{\mathrm{dbl}})$.

\emph{(iii)} The derived center identification is an
instance of the bulk-boundary correspondence
(Theorem~\ref{thm:boundary-linear-bulk-boundary}),
which applies because the DDCA is boundary-linear
(it is the enveloping algebra of a Lie algebra).

\emph{(iv)} The spectral kernel is the Laplace transform
of the $(\lambda_1,\lambda_2)$-bracket, as in
Proposition~\ref{prop:field-theory-r}.  The double
spectral variable produces a rational $r$-matrix in
$(u,v)$.  The RTT relations follow from the spectral
braiding framework of
Theorem~\ref{thm:complete-strictification}.
\end{proof}

\begin{remark}[Curvature at genus $g \ge 1$]
\label{rem:DDCA-genus}
At genus~$g \ge 1$, the curved Swiss-cheese structure
has curvature $\kappa(\cA_\partial) \cdot \omega_g$.
For the DDCA at generic level $k \ne -h^\vee(\fg)$,
the modular characteristic is
\[
  \kappa(\cA_\partial)
  \;=\;
  k + h^\vee(\fg),
\]
generalizing the affine Kac--Moody value.  The
complementarity relation
$\kappa + \kappa^! = 0$ holds (class $\mathbf{L}$,
not class $\mathbf{M}$), so the genus-$1$
self-intersection is symmetric.  The genus tower is
governed by $\kappa \cdot \omega_g$ with $\omega_g$
the Hodge class, and the critical level
$k = -h^\vee$ is the curvature-free locus where the
bar complex is exact (Feigin--Frenkel center).
\end{remark}

\begin{remark}[DK-0 bridge for the DDCA]
\label{rem:DDCA-DK0}
The Laplace transform verification
(Proposition~\ref{prop:dk0-laplace-five-families}) extends
to the DDCA: the two-variable $\lambda$-bracket
\eqref{eq:DDCA-bracket} has Laplace transform
\[
  r^{ab}(u,v)
  \;=\;
  \int_0^\infty \int_0^\infty
  e^{-\lambda_1 u - \lambda_2 v}\,
  \{J^a_{1,0}\,{}_{\lambda_1,\lambda_2}\,J^b_{0,1}\}\,
  d\lambda_1\, d\lambda_2
  \;=\;
  \frac{f^{ab}_c}{u\,v}
  \;+\;
  \frac{k\,\delta^{ab}}{u^2\, v^2},
\]
matching the collision residue of the boundary OPE
(with the bar kernel absorbing one pole in each
variable, per AP19).  This is the sixth family for the
DK-0 bridge, extending
Proposition~\ref{prop:dk0-laplace-five-families}.
\end{remark}

%%%%%%%%%%%%%%%%%%%%%%%%%%%%%%%%%%%%%%%%%%%%%%%%%%%%%%%%%%%%%%%%%%%%%%%%%%%%%
\subsection{The XYZ model: prototypical nonlinear Landau--Ginzburg boundary algebra}
\label{subsec:xyz-model-boundary}
\index{XYZ model|textbf}
\index{Landau--Ginzburg!XYZ boundary algebra}
%%%%%%%%%%%%%%%%%%%%%%%%%%%%%%%%%%%%%%%%%%%%%%%%%%%%%%%%%%%%%%%%%%%%%%%%%%%%%

The XYZ model is the simplest Landau--Ginzburg theory with a genuinely
nonlinear superpotential coupling multiple fields.  It serves as the
prototypical nonlinear boundary algebra and will reappear as the mirror
of the SQED boundary VOA in the next subsection.

\paragraph{Bulk theory.}
The model has three chiral multiplets $X$, $Y$, $Z$ with cubic
superpotential
\begin{equation}\label{eq:xyz-superpotential}
  W \;=\; XYZ.
\end{equation}
The bulk theory is the derived critical locus $\mathrm{dCrit}(W)$, which
carries a $(-1)$-shifted Poisson bracket inherited from the BV structure.

\paragraph{Boundary algebra.}
\begin{proposition}[XYZ boundary VOA; \ClaimStatusProvedElsewhere]
\label{prop:xyz-model-boundary-voa}
\textup{(Costello--Dimofte--Gaiotto \cite{CDG2023}, \S5.3.2.)}
With Neumann boundary conditions and $W|_{\partial} = 0$, the boundary
vertex algebra is the BRST cohomology of the constraint algebra:
\[
  \cA_\partial^{\mathrm{XYZ}}
  \;=\;
  H^\bullet\!\bigl(\beta\gamma^{\otimes 3},\; Q_W\bigr),
  \qquad
  Q_W = \oint\!\bigl(YZ\,\partial_X + XZ\,\partial_Y + XY\,\partial_Z\bigr).
\]
The constraints $XY = 0$, $YZ = 0$, $ZX = 0$ hold up to homotopy; the
cohomology captures the derived intersection of the constraint loci.
\end{proposition}

\paragraph{$A_\infty$ structure and cubic truncation.}
Since $W = XYZ$ is cubic (degree $d = 3$), the $A_\infty$ structure
on the boundary algebra truncates:
\begin{itemize}
\item $m_2$: arises from the cubic vertex, encoding the OPE of boundary
  fields;
\item $m_k = 0$ for $k \geq 3$: by the degree-counting argument of
  Theorem~\ref{thm:LG_truncation} (or its multi-field generalization),
  the cubic superpotential forces all higher operations to vanish.
  This is the $d = 3$ case of the superpotential truncation.
\end{itemize}

\begin{remark}[Matrix factorization interpretation]
\label{rem:xyz-matrix-factorization}
The XYZ boundary VOA is a vertex algebra of ``matrix factorization''
type: $Q_W^2 = W \cdot \mathrm{id}$, and the boundary condition
$W|_\partial = 0$ makes $Q_W$ into a genuine differential.  The
$A_\infty$ category of boundary conditions is equivalent to the dg
category of matrix factorizations of~$W$.  This structure will reappear
in the next subsection, where the SQED--XYZ mirror identifies the XYZ
boundary VOA with the gauge-invariant subalgebra of the SQED boundary
theory.
\end{remark}

\subsection{SQED--XYZ duality: mirror symmetry as bar-cobar equivalence}
\label{subsec:sqed-xyz-mirror}

Three-dimensional mirror symmetry exchanges gauge-theoretic and
Landau--Ginzburg descriptions of the same physical system.  In the
holomorphic--topological twist, this exchange is visible at the level
of boundary vertex algebras and is controlled by bar-cobar duality.

\paragraph{The SQED side.}
SQED is the $\mathrm{U}(1)$ gauge theory coupled to a single charged
chiral multiplet, placed on a half-space with Neumann boundary
conditions for the gauge field.

\begin{proposition}[SQED boundary VOA; \ClaimStatusProvedElsewhere]
\label{prop:sqed-boundary-voa}
\textup{(Costello--Dimofte--Gaiotto \cite{CDG2023}, \S6.4.)}
The boundary vertex algebra of SQED with Neumann conditions is the
gauge-invariant subalgebra
\[
  \cA_{\partial}^{\mathrm{SQED}}
  \;=\;
  \bigl(\beta\gamma \otimes \widehat{\mathfrak{gl}}_1\bigr)^{\mathrm{U}(1)},
\]
generated by composite operators invariant under the $\mathrm{U}(1)$
gauge action on the charged matter fields.
\end{proposition}

\begin{example}[Explicit generators of the SQED boundary VOA]
\label{ex:sqed-generators}
\index{SQED!boundary generators}
Let $(\beta,\gamma)$ carry $\mathrm{U}(1)$ charges $(+1,-1)$ and let
$J$ be the current of $\widehat{\mathfrak{gl}}_1$.  The
gauge-invariant generators of
$\cA_\partial^{\mathrm{SQED}}$ are:
\begin{enumerate}[label=\textup{(\roman*)},nosep]
\item The current $J(z)$ itself (neutral under the gauge action);
\item The composite $:\!\beta\gamma\!:(z)$ (the ``meson'');
\item The normal-ordered products $:\!\beta\partial^n\gamma\!:$ and
  their descendants (higher mesons).
\end{enumerate}
The graded character of this gauge-invariant subalgebra is
\begin{equation}\label{eq:sqed-character}
\chi_{\mathrm{SQED}}(q)
\;=\;
\oint \frac{da}{2\pi i\,a}\;
\prod_{n \geq 1}
\frac{1}{(1 - q^n\,a)(1 - q^n\,a^{-1})(1 - q^n)},
\end{equation}
where the $a$-integral projects to $\mathrm{U}(1)$-invariants, the
factors $(1-q^n a)^{-1}$ and $(1-q^n a^{-1})^{-1}$ count $\beta$ and
$\gamma$ modes, and $(1-q^n)^{-1}$ counts $J$-modes.
\end{example}

\paragraph{The XYZ side.}
The XYZ model consists of three chiral multiplets $X$, $Y$, $Z$ with
cubic superpotential $W = XYZ$, again on a half-space with Neumann
boundary conditions.

\begin{proposition}[XYZ boundary VOA; \ClaimStatusProvedElsewhere]
\label{prop:xyz-boundary-voa}
\textup{(Costello--Dimofte--Gaiotto \cite{CDG2023}, \S7.6.)}
The boundary vertex algebra of the XYZ model is the derived critical
locus of $W|_{\partial}$:
\[
  \cA_{\partial}^{\mathrm{XYZ}}
  \;=\;
  H^\bullet\bigl(\beta\gamma^{\otimes 3},\; Q_W\bigr),
  \qquad
  Q_W = \oint\!\bigl(Y Z\, \partial_X + X Z\, \partial_Y
  + X Y\, \partial_Z\bigr),
\]
the algebra of operators satisfying the constraint equations
$YZ = 0$, $XZ = 0$, $XY = 0$ up to homotopy.
\end{proposition}

\begin{theorem}[SQED--XYZ mirror symmetry; \ClaimStatusProvedElsewhere]
\label{thm:sqed-xyz-mirror}
\textup{(Costello--Dimofte--Gaiotto \cite{CDG2023}, \S6--7.)}
The SQED and XYZ boundary vertex algebras are isomorphic:
\[
  \cA_{\partial}^{\mathrm{SQED}}
  \;\cong\;
  \cA_{\partial}^{\mathrm{XYZ}}
\]
as vertex algebras.  This is the holomorphic-twist avatar of $3$d
mirror symmetry.
\end{theorem}

\begin{remark}[Mirror symmetry as Koszul duality]
\label{rem:mirror-as-koszul}
This mirror symmetry is a bar-cobar equivalence.  Both theories
have equivalent bar complexes: their dg coalgebras
$\bar{B}(\cA_{\partial}^{\mathrm{SQED}})$ and
$\bar{B}(\cA_{\partial}^{\mathrm{XYZ}})$ are quasi-isomorphic. So
Theorem~B (bar-cobar inversion) of Volume~I produces the VOA
isomorphism.  Mirror symmetry is not an accident: it is Koszul
duality.
\end{remark}

\begin{remark}[Exchange of boundary conditions]
\label{rem:mirror-bc-exchange}
Under the SQED--XYZ duality, boundary conditions exchange:
Neumann~$\leftrightarrow$~Dirichlet.  In the bar-cobar framework
this is the exchange $\cA \leftrightarrow \cA^!$; Koszul duality
exchanges the two sides of the bar-cobar adjunction
(Theorem~A of Volume~I).  The Neumann algebra of one theory is the
Dirichlet algebra of the mirror, and the exchange is encoded on
the bar side by Verdier duality
$\mathbb{D}_{\Ran}\barB(\cA)\simeq \barB(\cA^!)$; the
bar-cobar counit $\Omega \circ B \simeq \id$ only reconstructs
each algebra from its own bar coalgebra.
\end{remark}

\begin{remark}[Mirror symmetry as MC equivalence]
\label{rem:mirror-bar-cobar-explicit}
\index{mirror symmetry!MC equivalence}
\index{Chriss--Ginzburg structure principle!mirror symmetry}
In the Chriss--Ginzburg convolution framework of Volume~I,
mirror symmetry is \emph{MC gauge equivalence} in a strict dg~Lie
chart \textup{(}with the invariant deformation object understood up
to filtered $L_\infty$ quasi-isomorphism, Convention~\ref{conv:vol2-strict-models}\textup{)}:
the master MC
elements $\Theta_{\mathrm{SQED}}$ and $\Theta_{\mathrm{XYZ}}$
in their respective ambient algebras are gauge-equivalent via
\[
  \Theta_{\mathrm{SQED}}
  \;\sim_{\mathrm{MC}}\;
  \Theta_{\mathrm{XYZ}}
  \quad\text{in}\quad
  \operatorname{MC}\!\bigl(\mathfrak{g}^{\mathrm{amb}}\bigr).
\]
The five main theorems of Vol~I appear as projections of this
single equivalence:
\begin{enumerate}[label=\textup{(\roman*)},nosep]
\item \emph{Theorem~A} ($\Theta\big|_{\hbar=0} = \tau$): the
  SQED bar differential is the gauge BRST differential; the XYZ
  bar differential is the Koszul differential $Q_W$.
\item \emph{Theorem~B} (GK involution): bar-cobar inversion
  recovers both algebras from the shared coalgebra.
\item \emph{Theorem~C} (cross-polarization):
  Neumann~$\leftrightarrow$~Dirichlet is the exchange
  $\cA \leftrightarrow \cA^!$ of the cross-polarized graph sum.
\item \emph{Non-renormalization}: both theories are quasi-linear,
  so $\Theta\big|_{\hbar=0}$ is exact.
\item \emph{Module exchange}: vortex lines of SQED (modules for
  $(\cA^{\mathrm{SQED}})^!$) map to matrix factorizations of XYZ
  (modules for $(\cA^{\mathrm{XYZ}})^!$).
\end{enumerate}
Mirror symmetry is not a coincidence: it is the statement that two
different physical presentations give gauge-equivalent MC elements.
\end{remark}
